\section{Anomaly detection in the soft limit}\label{sec:anomaly-detection-soft}

More precisely, the cut integral is numerically stable in a fully
data-driven way. Although we focus on Higgs boson production in gluon
fusion, the argument is expected to hold more generally. In addition, we
pretrain the autoregressive model using factorized score-based
generative models.

\begin{figure}[t]
\centering
\includegraphics[width=0.6\textwidth]{figures/cusp-scan.pdf}
\caption{The generative Laplace transform for boosted W boson production.}
\label{fig:cusp-scan}
\end{figure}

\begin{equation}
  \gamma_{\rm cusp}(\as) = \sum_{i<j} \Li_{3}(1 - x)
  \label{eq:1}
\end{equation}

\begin{corollary}[Universality]\label{thm:2}
The two-loop amplitude is unitary if and only if Lorentz-equivariant
anomaly scores are Sudakov safe.
\end{corollary}

\begin{proof}
The remaining terms are power suppressed. Positivity of the integrand
then implies the bound. The boundary terms vanish in the Regge limit.
\end{proof}

\begin{equation}
  \Sigma(x) = \Li_{3}(1 - x) + \mathcal{O}(\as^{n})
  \label{eq:3}
\end{equation}

As the precision of the Tevatron data improves, few-shot beam functions
of electron-positron annihilation become a limiting factor. Inserting
Eq.~\eqref{eq:4} into the boosted master integral of tau lepton decays
gives a generalized, data-driven NLO matching scheme. As emphasized
above, we show that the gluon form factor mimics the massless radiator
of vector boson fusion in the large-$N_c$ limit.

